\chapter{The Four Fundamental Subspaces}

You have now met two subspaces that live inside every matrix: the column space, where the
outputs $\mat{A}\vec{x}$ live, and the nullspace, where the inputs that get crushed to zero
live. In this chapter we meet their two partners, and we discover something wonderful: a single
$m \times n$ matrix quietly organizes \emph{four} subspaces -- two sitting in the input space
$\R^n$, two sitting in the output space $\R^m$ -- and these four spaces fit together with such
geometric precision that Gilbert Strang calls the resulting figure ``the big picture of linear
algebra.'' Once you see the picture, you understand what a matrix really does: it takes one
$r$-dimensional subspace and maps it cleanly, invertibly, onto another $r$-dimensional subspace,
while flattening everything else to nothing.

That is the whole story. Everything below is us slowing down to see it from the ground up.

\section{Two spaces for the inputs, two for the outputs}

Let $\mat{A}$ be an $m \times n$ matrix. It eats vectors $\vec{x} \in \R^n$ and produces
vectors $\mat{A}\vec{x} \in \R^m$. So already there are two ``rooms'' in the building: the
\emph{input room} $\R^n$ and the \emph{output room} $\R^m$. Each room will turn out to contain
exactly two of our fundamental subspaces, and they will be \emph{orthogonal complements} of
each other inside that room.

\begin{focusbox}
The four fundamental subspaces of an $m \times n$ matrix $\mat{A}$ are:
\begin{itemize}
  \item the \textbf{column space} $\col(\mat{A}) \subseteq \R^m$ -- all outputs $\mat{A}\vec{x}$;
  \item the \textbf{nullspace} $\nullsp(\mat{A}) \subseteq \R^n$ -- all inputs killed by $\mat{A}$;
  \item the \textbf{row space} $\row(\mat{A}) = \col(\mat{A}\T) \subseteq \R^n$ -- the span of the rows;
  \item the \textbf{left nullspace} $\nullsp(\mat{A}\T) \subseteq \R^m$ -- inputs killed by $\mat{A}\T$.
\end{itemize}
Two live in $\R^n$ (column-vector inputs): the row space and the nullspace.
Two live in $\R^m$ (column-vector outputs): the column space and the left nullspace.
\end{focusbox}

\begin{intuitionbox}
Picture a matrix as a machine that pours $\R^n$ into $\R^m$. Inside $\R^n$, the machine secretly
splits every input into two pieces: a piece it \emph{listens to} (the row space) and a piece it
\emph{ignores} (the nullspace). Inside $\R^m$, the outputs only ever land in a flat slab (the
column space); the rest of $\R^m$ (the left nullspace) is unreachable. The magic is that the
machine maps the ``listened-to'' part of the input perfectly onto the ``reachable'' part of the
output, one-to-one, no information lost.
\end{intuitionbox}

Before the big picture, let us build each space carefully and learn to read its dimension and a
basis straight off the elimination of $\mat{A}$.

\section{The column space: where the outputs live}

Recall the definition. The column space is everything you can build out of the columns of
$\mat{A}$:
\[
  \col(\mat{A}) \;=\; \set{\,\mat{A}\vec{x} : \vec{x}\in\R^n\,}
  \;=\; \spn(\text{columns of }\mat{A}) \;\subseteq\; \R^m.
\]
The equation $\mat{A}\vec{x}=\vec{b}$ has a solution exactly when $\vec{b}$ lies in this space.
The column space is the set of \emph{achievable right-hand sides}.

\begin{example}\label{ch15:colspace-rank2}
Take
\[
  \mat{A} \;=\;
  \begin{bmatrix} 1 & 2 & 3\\ 2 & 4 & 8\\ 1 & 2 & 5 \end{bmatrix}.
\]
Column~2 is exactly $2\times$ column~1, so it adds nothing new. Columns~1 and~3 point in genuinely
different directions. So the column space is the plane through the origin in $\R^3$ spanned by
$\colvec{1\\2\\1}$ and $\colvec{3\\8\\5}$. It is two-dimensional: a flat sheet inside three-space.
The vector $\colvec{0\\0\\1}$, say, is \emph{not} on that sheet, so $\mat{A}\vec{x}=\colvec{0\\0\\1}$
has no solution.
\end{example}

How did we know columns~1 and~3 were the ``real'' ones? Elimination tells us. The pivot columns
of $\mat{A}$ -- the columns that end up holding a pivot after row reduction -- form a basis for the
column space.

\begin{formalbox}
\textbf{Basis for the column space.} Row-reduce $\mat{A}$ to find which columns contain pivots.
The corresponding columns \emph{of the original $\mat{A}$} form a basis for $\col(\mat{A})$.
Therefore $\dimn \col(\mat{A}) = r$, the number of pivots, which we call the \textbf{rank}.
\end{formalbox}

\begin{trapbox}
A subtle but classic mistake: after row reduction you must take the pivot columns of the
\emph{original} matrix, not the pivot columns of the reduced matrix. Row operations change the
column space! (They mix the rows, which moves the columns around in $\R^m$.) What row operations
\emph{do} preserve is \emph{which} columns are dependent on which -- so the \emph{positions} of the
pivots are trustworthy, but the actual basis vectors must come from $\mat{A}$ itself.
\end{trapbox}

Let us see the trap in miniature.

\begin{labbox}
Reduce $\mat{A}=\begin{bmatrix}1&1\\1&1\end{bmatrix}$. You get
$\begin{bmatrix}1&1\\0&0\end{bmatrix}$. The pivot is in column~1. The column space of the reduced
matrix is the span of $\colvec{1\\0}$ (the $x$-axis). But the column space of $\mat{A}$ is the
span of $\colvec{1\\1}$ (the diagonal line). Different lines! Same dimension ($r=1$), same pivot
positions, but you must read the basis vector $\colvec{1\\1}$ off the original $\mat{A}$.
\end{labbox}

\begin{figure}[H]
\centering
\begin{tikzpicture}[scale=1.1,>=Stealth]
  % axes
  \draw[->,spacegray] (-0.3,0)--(3.2,0) node[right]{$x$};
  \draw[->,spacegray] (0,-0.3)--(0,3.2) node[above]{$y$};
  % column space of A: the diagonal
  \draw[vecblue,very thick] (-0.2,-0.2)--(3,3);
  \fill[vecblue] (1,1) circle (1.6pt);
  \draw[->,vecblue,very thick] (0,0)--(1,1);
  \node[vecblue,right] at (1.05,0.85){$\col(\mat{A})$};
  % column space of reduced: x-axis (dashed red)
  \draw[vecred,very thick,dashed] (-0.2,0)--(3,0);
  \draw[->,vecred,very thick] (0,0)--(1,0);
  \node[vecred,below] at (2.2,-0.05){$\col(\text{reduced})$};
\end{tikzpicture}
\caption{Row reduction preserves pivot \emph{positions} but not the column space itself. For
$\mat{A}=\protect\begin{bmatrix}1&1\\1&1\protect\end{bmatrix}$ the true column space is the
diagonal line (blue); the reduced matrix has the $x$-axis (red dashed). Always read the basis
off the original columns.}
\label{ch15:colspace-trap}
\end{figure}

\section{The nullspace: where the inputs die}

The nullspace collects every input that $\mat{A}$ sends to zero:
\[
  \nullsp(\mat{A}) \;=\; \set{\,\vec{x}\in\R^n : \mat{A}\vec{x}=\vec{0}\,} \;\subseteq\; \R^n.
\]
This is a subspace because if $\mat{A}\vec{x}=\vec{0}$ and $\mat{A}\vec{y}=\vec{0}$ then
$\mat{A}(c\vec{x}+d\vec{y})=\vec{0}$ too. Geometrically it is the set of directions the machine
is blind to -- push along a nullspace direction and the output does not budge.

We find a basis by solving $\mat{A}\vec{x}=\vec{0}$ with elimination. The \emph{free} variables
(columns with no pivot) are the dials we get to turn; each one, set to $1$ while the others are
$0$, produces one \emph{special solution}, and the special solutions form a basis.

\begin{example}\label{ch15:nullspace-special}
Continue with $\mat{A}=\begin{bmatrix}1&2&3\\2&4&8\\1&2&5\end{bmatrix}$ from
Example~\ref{ch15:colspace-rank2}. Eliminate:
\[
  \begin{bmatrix}1&2&3\\2&4&8\\1&2&5\end{bmatrix}
  \;\xrightarrow{R_2-2R_1,\,R_3-R_1}\;
  \begin{bmatrix}1&2&3\\0&0&2\\0&0&2\end{bmatrix}
  \;\xrightarrow{R_3-R_2}\;
  \begin{bmatrix}1&2&3\\0&0&2\\0&0&0\end{bmatrix}.
\]
Pivots sit in columns~1 and~3; column~2 is free. Set the free variable $x_2 = 1$. Back-substitute:
the second row gives $2x_3=0$, so $x_3=0$; the first row gives $x_1+2x_2+3x_3=0$, so $x_1=-2$. The
special solution is
\[
  \vec{s} = \colvec{-2\\1\\0},
  \qquad
  \nullsp(\mat{A}) = \spn\set[\big]{\colvec{-2\\1\\0}}.
\]
Check: $\mat{A}\vec{s}=\colvec{1\cdot(-2)+2\cdot1+3\cdot0\\ 2\cdot(-2)+4\cdot1+8\cdot0\\
1\cdot(-2)+2\cdot1+5\cdot0}=\colvec{0\\0\\0}.$ One free variable, so the nullspace is a single
line in $\R^3$: $\dimn\nullsp(\mat{A})=1$.
\end{example}

Notice the bookkeeping that just happened. There were $n=3$ columns. Two held pivots ($r=2$), one
was free ($n-r=1$). That free variable became the one dimension of the nullspace. This is no
accident.

\begin{formalbox}
\textbf{Dimension of the nullspace.} For an $m\times n$ matrix of rank $r$, the number of free
variables is $n-r$, and each gives an independent special solution. Hence
\[
  \dimn \nullsp(\mat{A}) \;=\; n - r.
\]
This is the \emph{nullity} of $\mat{A}$, and the identity $\rank + \text{nullity} = n$ is the
celebrated \textbf{rank--nullity theorem}.
\end{formalbox}

\begin{walkbox}
Take a walk and chew on this: the rank counts the directions the matrix genuinely \emph{uses},
the nullity counts the directions it \emph{wastes}, and the two must add up to all $n$ input
directions. There is nowhere else for a direction to go -- it is either heard or ignored. Why does
that feel inevitable once you say it out loud?
\end{walkbox}

\section{The row space: the input directions the matrix listens to}

The rows of $\mat{A}$ are themselves vectors -- but they have $n$ components, so they live in the
\emph{input} room $\R^n$. Their span is the row space:
\[
  \row(\mat{A}) \;=\; \spn(\text{rows of }\mat{A}) \;=\; \col(\mat{A}\T) \;\subseteq\; \R^n.
\]
Why call this ``the directions the matrix listens to''? Because the $i$-th component of the output
is $\mat{A}\vec{x}$ read row by row: $(\mat{A}\vec{x})_i = (\text{row }i)\cdot \vec{x}$. The output
is built entirely from dot products of $\vec{x}$ against the rows. If $\vec{x}$ is orthogonal to
every row, every dot product is zero and the output vanishes -- that is exactly the nullspace. So
the row space and the nullspace are perpendicular partners inside $\R^n$. Hold that thought; we
make it precise in the next section.

How big is the row space, and what is a basis? Here the news is friendlier than for the column
space.

\begin{formalbox}
\textbf{Basis for the row space.} Row operations are reversible combinations of rows, so they
never change the span of the rows. Therefore the nonzero rows of the \emph{row-reduced} matrix
$\mat{R}$ form a basis for $\row(\mat{A})=\row(\mat{R})$, and
\[
  \dimn \row(\mat{A}) \;=\; r \;=\; \dimn\col(\mat{A}).
\]
Unlike the column space, you \emph{may} use the rows of the reduced matrix directly.
\end{formalbox}

That last equation is quietly astonishing: the number of independent rows equals the number of
independent columns. We give it the attention it deserves in the next chapter, but you have already
proved it -- both counts equal $r$, the number of pivots.

\begin{example}\label{ch15:rowspace}
Same $\mat{A}$. Its reduced rows (from Example~\ref{ch15:nullspace-special}, after scaling the
pivot to $1$) are
\[
  \row(\mat{A}) = \spn\set[\Big]{\colvec{1\\2\\3},\ \colvec{0\\0\\1}}.
\]
A two-dimensional plane in $\R^3$. And look: the nullspace direction $\colvec{-2\\1\\0}$ is
orthogonal to both row vectors. Dot with $\colvec{1\\2\\3}$: $-2+2+0=0$. Dot with
$\colvec{0\\0\\1}$: $0+0+0=0$. The plane (row space) and the line (nullspace) meet at a perfect
right angle, and together $2+1=3$ they fill all of $\R^3$.
\end{example}

Figure~\ref{ch15:rowspace-null} shows this split of the input room: the row space is a plane, the
nullspace is the line through it at a right angle, and every input vector decomposes uniquely into a
piece in each.

\begin{figure}[H]
\centering
\begin{tikzpicture}[scale=1.0,>=Stealth]
  % a tilted plane (row space) drawn as a parallelogram, with a normal line (nullspace)
  \fill[sourceblue!14,draw=sourceblue!70!black,thick]
        (-2.1,-0.9) -- (1.9,-1.5) -- (2.1,0.9) -- (-1.9,1.5) -- cycle;
  \node[sourceblue!70!black] at (-1.35,1.05) {$\row(\mat{A})$ (plane, $\dim 2$)};
  % nullspace line, perpendicular, drawn rising out of the plane
  \draw[quotientorange!85!black,very thick] (0,0) -- (0.9,2.3);
  \draw[quotientorange!85!black,very thick,dashed] (0,0) -- (-0.9,-2.3);
  \draw[->,quotientorange!85!black,very thick] (0,0) -- (0.7,1.8);
  \node[quotientorange!85!black,right] at (0.7,1.95) {$\nullsp(\mat{A})$ (line, $\dim 1$)};
  % right-angle hint at origin
  \node[fill=black,circle,inner sep=1.1pt] at (0,0) {};
  \node[font=\scriptsize,spacegray,below right] at (0.05,-0.05) {$\vec 0$};
  \node[font=\scriptsize,spacegray] at (0.42,0.4) {$\perp$};
\end{tikzpicture}
\caption{The input room $\R^3$ for the running matrix
(Example~\ref{ch15:rowspace}), split into perpendicular partners: the row space is the plane
spanned by $\protect\colvec{1\\2\\3}$ and $\protect\colvec{0\\0\\1}$, and the nullspace is the line
through $\protect\colvec{-2\\1\\0}$, which stands at a right angle to that plane. Dimensions
$2+1=3$ fill the room.}
\label{ch15:rowspace-null}
\end{figure}

\begin{strangbox}
\emph{Before Strang:} when you first watch him reduce a matrix to find ``pivot columns'' and
``free columns,'' it can look like mechanical bookkeeping. Carry this chapter's picture with you
and it becomes geometry: pivots count the dimensions of the row and column spaces; free columns
count the dimensions of the nullspace.

\emph{After Strang:} he draws the four-subspace figure on the board more than any other diagram
in his course and calls it the heart of the subject. If you have read this chapter, you will
recognize it instantly and you will know \emph{why} the two spaces in each room are perpendicular.
\end{strangbox}

\section{The left nullspace: the unreachable output directions}

The fourth space is the most easily forgotten because we have to transpose to name it. The
\emph{left nullspace} is the nullspace of $\mat{A}\T$:
\[
  \nullsp(\mat{A}\T) \;=\; \set{\,\vec{y}\in\R^m : \mat{A}\T\vec{y}=\vec{0}\,} \;\subseteq\; \R^m.
\]
Why ``left''? Because $\mat{A}\T\vec{y}=\vec{0}$ is the same as $\vec{y}\T\mat{A}=\vec{0}\T$ -- a
row vector $\vec{y}\T$ multiplying $\mat{A}$ from the \emph{left} and producing the zero row. So
$\vec{y}$ records a linear combination of the \emph{rows} of $\mat{A}$ that cancels to zero. These
are exactly the redundancies among the rows.

Geometrically, the left nullspace lives in the output room $\R^m$ and consists of the directions
the column space cannot reach. If $\vec{y}\in\nullsp(\mat{A}\T)$, then for every input $\vec{x}$,
\[
  \vec{y}\cdot(\mat{A}\vec{x}) = \vec{y}\T\mat{A}\vec{x} = (\mat{A}\T\vec{y})\T\vec{x}
  = \vec{0}\T\vec{x}=0.
\]
So $\vec{y}$ is perpendicular to every possible output. The left nullspace is the orthogonal
complement of the column space inside $\R^m$.

\begin{formalbox}
\textbf{Dimension and basis of the left nullspace.} Applying the rank--nullity bookkeeping to
$\mat{A}\T$ (which is $n\times m$ with the same rank $r$):
\[
  \dimn \nullsp(\mat{A}\T) \;=\; m - r.
\]
A clean way to get a basis: do elimination as $\mat{E}\mat{A}=\mat{R}$, tracking the operations
in $\mat{E}$. The rows of $\mat{E}$ that produce the \emph{zero rows} of $\mat{R}$ are exactly a
basis for the left nullspace.
\end{formalbox}

\begin{example}\label{ch15:leftnull}
For our running $\mat{A}$, the reduction $R_2-2R_1$, $R_3-R_1$, then $R_3-R_2$ combined the rows
so that the final third row was all zeros. Tracking the combination, the third row of $\mat{R}$
came from
\[
  (R_3 - R_1) - (R_2 - 2R_1) = R_1 - R_2 + R_3,
\]
i.e. row~1 minus row~2 plus row~3 of $\mat{A}$ gives the zero row. So
$\vec{y}=\colvec{1\\-1\\1}$ satisfies $\vec{y}\T\mat{A}=\vec{0}\T$. Check the first component:
$1\cdot1 - 1\cdot2 + 1\cdot1 = 0$. Second: $1\cdot2-1\cdot4+1\cdot2=0$. Third:
$1\cdot3-1\cdot8+1\cdot5=0$. Indeed $\vec{y}\T\mat{A}=\vec{0}\T$. Here $m=3$, $r=2$, so the left
nullspace is the single line $\spn\set{\colvec{1\\-1\\1}}$ in $\R^3$, dimension $m-r=1$. And it is
orthogonal to the column space plane: dot $\colvec{1\\-1\\1}$ with column~1 $\colvec{1\\2\\1}$ to
get $1-2+1=0$, and with column~3 $\colvec{3\\8\\5}$ to get $3-8+5=0$. Right angles everywhere.
\end{example}

\begin{aestheticbox}
Step back and admire the symmetry. The matrix and its transpose share the same rank $r$. So all
four dimensions are determined by just three numbers $m$, $n$, $r$:
\[
  \dimn\col=r,\quad \dimn\nullsp(\mat{A}\T)=m-r \quad(\text{in }\R^m),\qquad
  \dimn\row=r,\quad \dimn\nullsp=n-r \quad(\text{in }\R^n).
\]
In each room the two dimensions sum to the dimension of the room: $r+(m-r)=m$ and $r+(n-r)=n$.
The whole input space and the whole output space are each split cleanly into a ``used'' part and
an ``unused'' part. That a single rank number governs all of this is one of the quiet miracles of
the subject.
\end{aestheticbox}

\section{Orthogonality: the two spaces in each room are perpendicular}

We keep saying ``perpendicular partners.'' Let us prove it cleanly, because it is the structural
backbone of the big picture.

\begin{theorem}[Orthogonality of the fundamental subspaces]\label{ch15:ortho-thm}
For any $m\times n$ matrix $\mat{A}$:
\begin{enumerate}
  \item In $\R^n$: every nullspace vector is orthogonal to every row-space vector. In symbols,
  $\nullsp(\mat{A}) \perp \row(\mat{A})$, and together they fill $\R^n$.
  \item In $\R^m$: every left-nullspace vector is orthogonal to every column-space vector. In
  symbols, $\nullsp(\mat{A}\T) \perp \col(\mat{A})$, and together they fill $\R^m$.
\end{enumerate}
\end{theorem}

\begin{proof}
\emph{Part 1.} Let $\vec{x}\in\nullsp(\mat{A})$, so $\mat{A}\vec{x}=\vec{0}$. Reading $\mat{A}\vec{x}$
row by row, the $i$-th entry is $(\text{row}_i)\cdot\vec{x}=0$. So $\vec{x}$ is orthogonal to every
row, hence to every linear combination of rows, hence to the entire row space. That gives
$\nullsp(\mat{A})\perp\row(\mat{A})$. For the dimension count: $\dimn\row=r$ and $\dimn\nullsp=n-r$,
and since the spaces meet only at $\vec{0}$ (a nonzero shared vector would be orthogonal to itself,
forcing length zero), their dimensions add to $r+(n-r)=n=\dimn\R^n$. Two orthogonal subspaces whose
dimensions fill the room must \emph{be} orthogonal complements: every vector in $\R^n$ splits
uniquely into a row-space piece plus a nullspace piece.

\emph{Part 2} is Part 1 applied to $\mat{A}\T$: its row space is $\col(\mat{A})$ and its nullspace
is $\nullsp(\mat{A}\T)$, both living in $\R^m$.
\end{proof}

\begin{intuitionbox}
Here is the cleanest way to remember which pairs are perpendicular. The nullspace is defined by
``dot with every row gives zero,'' so the nullspace is whatever is left after you remove the row
space -- they are complements. Likewise the left nullspace is ``dot with every column gives zero,''
so it is the complement of the column space. Rows pair with the nullspace; columns pair with the
left nullspace. The output ($\col$) and the unreachable directions ($\nullsp(\mat{A}\T)$) carve up
$\R^m$; the listened-to inputs ($\row$) and the ignored inputs ($\nullsp$) carve up $\R^n$.
\end{intuitionbox}

\begin{connectionbox}
The phrase ``orthogonal complement and the spaces fill the room'' is the
\emph{Fundamental Theorem of Linear Algebra, Part 2}, and it is the finite-dimensional shadow of a
deep theme in analysis. In the study of differential and integral operators (functional analysis),
the same statement is the \emph{Fredholm alternative}: the equation $\mat{A}\vec{x}=\vec{b}$ is
solvable precisely when $\vec{b}$ is orthogonal to every solution of the ``adjoint'' homogeneous
problem $\mat{A}\T\vec{y}=\vec{0}$. That adjoint condition is exactly ``$\vec{b}\perp$ left
nullspace, i.e. $\vec{b}\in\col(\mat{A})$.'' The four-subspace picture you are learning here scales
all the way up to infinite-dimensional operators and underwrites the solvability theory for huge
classes of differential equations.
\end{connectionbox}

\section{The big picture}

Now we assemble everything into the one diagram worth memorizing. Draw the input room $\R^n$ on
the left and the output room $\R^m$ on the right. Inside each room, draw the two perpendicular
subspaces. Then draw the action of $\mat{A}$ as arrows crossing from left to right.

\begin{figure}[H]
\centering
\begin{tikzpicture}[scale=1.0,>=Stealth,
   space/.style={draw,rounded corners,minimum width=2.6cm,minimum height=1.05cm,align=center,thick},
   row/.style={space,fill=sourceblue!12,draw=sourceblue},
   nul/.style={space,fill=quotientorange!16,draw=quotientorange!80!black},
   col/.style={space,fill=targetgreen!14,draw=targetgreen!80!black},
   lnul/.style={space,fill=structurepurple!14,draw=structurepurple}]

  % left room R^n
  \node[row]  (rs)  at (0,1.6)  {row space\\[-1pt]$\dim = r$};
  \node[nul]  (ns)  at (0,-1.6) {nullspace\\[-1pt]$\dim = n-r$};
  \node[spacegray] at (0,3.0) {\textbf{$\R^n$ (inputs)}};

  % right room R^m
  \node[col]  (cs)  at (6.4,1.6)  {column space\\[-1pt]$\dim = r$};
  \node[lnul] (lns) at (6.4,-1.6) {left nullspace\\[-1pt]$\dim = m-r$};
  \node[spacegray] at (6.4,3.0) {\textbf{$\R^m$ (outputs)}};

  % origins
  \node[circle,fill=black,inner sep=1.1pt] (o1) at (0,0) {};
  \node[circle,fill=black,inner sep=1.1pt] (o2) at (6.4,0) {};
  \node[font=\scriptsize] at (0.0,-0.32) {$\vec 0$};
  \node[font=\scriptsize] at (6.4,-0.32) {$\vec 0$};

  % the isomorphism row space -> column space
  \draw[->,very thick,sourceblue] (rs.east) to[bend left=12] node[above,font=\small]{$\mat{A}$: one-to-one, onto} (cs.west);
  % nullspace collapses to origin
  \draw[->,very thick,quotientorange!80!black] (ns.east) to[bend right=14] node[below,font=\small,pos=0.55]{$\mat{A}$: everything $\mapsto \vec 0$} (o2);

  % perpendicular ticks (just a hint)
  \node[font=\scriptsize,spacegray] at (0,0) {$\perp$};
  \node[font=\scriptsize,spacegray] at (6.4,0) {$\perp$};
\end{tikzpicture}
\caption{The big picture. The input room $\R^n$ splits into the row space and nullspace
(perpendicular, dimensions $r$ and $n-r$). The output room $\R^m$ splits into the column space and
left nullspace (perpendicular, dimensions $r$ and $m-r$). The matrix $\mat{A}$ collapses the
nullspace to $\vec 0$ and maps the row space \emph{one-to-one and onto} the column space.}
\label{ch15:bigpicture}
\end{figure}

Read the figure slowly. Two facts make it click.

First, \textbf{the nullspace collapses}. Any input in the nullspace goes straight to $\vec 0$ in
the output room. That is its definition.

Second -- and this is the beautiful part -- \textbf{the row space maps perfectly onto the column
space}. Take any input $\vec{x}\in\R^n$ and split it (uniquely, by
Theorem~\ref{ch15:ortho-thm}) as
\[
  \vec{x} = \vec{x}_{\text{row}} + \vec{x}_{\text{null}},
  \qquad \vec{x}_{\text{row}}\in\row(\mat{A}),\ \ \vec{x}_{\text{null}}\in\nullsp(\mat{A}).
\]
Then
\[
  \mat{A}\vec{x} = \mat{A}\vec{x}_{\text{row}} + \mat{A}\vec{x}_{\text{null}}
  = \mat{A}\vec{x}_{\text{row}} + \vec{0} = \mat{A}\vec{x}_{\text{row}}.
\]
So the output is entirely determined by the row-space part of the input -- the nullspace part is
silently discarded. The map ``$\vec{x}_{\text{row}}\mapsto\mat{A}\vec{x}_{\text{row}}$'' goes from
the row space (dimension $r$) to the column space (dimension $r$).

\begin{theorem}[Row space to column space is an isomorphism]\label{ch15:iso-thm}
Restricted to the row space, $\mat{A}$ is a one-to-one and onto linear map from $\row(\mat{A})$ to
$\col(\mat{A})$.
\end{theorem}

\begin{proof}
\emph{Onto:} every output $\mat{A}\vec{x}$ equals $\mat{A}\vec{x}_{\text{row}}$ by the splitting
above, so every element of the column space is hit by some row-space vector.

\emph{One-to-one:} suppose two row-space vectors $\vec{u},\vec{v}\in\row(\mat{A})$ give the same
output, $\mat{A}\vec{u}=\mat{A}\vec{v}$. Then $\mat{A}(\vec{u}-\vec{v})=\vec{0}$, so
$\vec{u}-\vec{v}\in\nullsp(\mat{A})$. But $\vec{u}-\vec{v}$ is also in the row space (a difference
of row-space vectors). A vector in both the row space and the nullspace is orthogonal to itself,
so $\vec{u}-\vec{v}=\vec{0}$, i.e. $\vec{u}=\vec{v}$. Distinct row-space vectors give distinct
outputs.
\end{proof}

\begin{focusbox}
This is the sentence to carry forever: \emph{every matrix is invertible if you look at it from the
right rooms}. Throw away the nullspace on the input side, throw away the left nullspace on the
output side, and what remains -- row space to column space -- is a perfect, reversible,
$r$-dimensional correspondence. The rank $r$ is the true ``size'' of the transformation.
\end{focusbox}

\begin{example}\label{ch15:splitting}
Let us watch the splitting in action with our $\mat{A}$. Take the input $\vec{x}=\colvec{0\\1\\0}$
(just turning on the free variable, plus nothing else). Decompose it. The nullspace direction is
$\vec{n}=\colvec{-2\\1\\0}$ with $\norm{\vec{n}}^2=5$. Project $\vec{x}$ onto it:
\[
  \vec{x}_{\text{null}} = \frac{\vec{x}\cdot\vec{n}}{\vec{n}\cdot\vec{n}}\,\vec{n}
  = \frac{(0)(-2)+(1)(1)+0}{5}\colvec{-2\\1\\0}
  = \frac{1}{5}\colvec{-2\\1\\0}.
\]
Then $\vec{x}_{\text{row}} = \vec{x}-\vec{x}_{\text{null}}
= \colvec{0\\1\\0}-\tfrac{1}{5}\colvec{-2\\1\\0} = \tfrac{1}{5}\colvec{2\\4\\0}$. You can verify
$\vec{x}_{\text{row}}$ lies in the row space (it is $\tfrac{2}{5}\colvec{1\\2\\0}$, a combination of
the rows) and is orthogonal to $\vec{n}$. Now apply $\mat{A}$:
\[
  \mat{A}\vec{x} = \mat{A}\colvec{0\\1\\0}=\colvec{2\\4\\2},\qquad
  \mat{A}\vec{x}_{\text{row}} = \mat{A}\cdot\tfrac15\colvec{2\\4\\0}
  =\tfrac15\colvec{2\cdot1+4\cdot2\\2\cdot2+4\cdot4\\2\cdot1+4\cdot2}=\tfrac15\colvec{10\\20\\10}
  =\colvec{2\\4\\2}.
\]
Identical. The nullspace part contributed nothing; the row-space part did all the work. That is
Theorem~\ref{ch15:iso-thm} with numbers.
\end{example}

Figure~\ref{ch15:split-arrows} pictures the mechanism: the input $\vec{x}$ casts a shadow onto the
row space and onto the nullspace; $\mat{A}$ keeps the row-space shadow and erases the other.

\begin{figure}[H]
\centering
\begin{tikzpicture}[scale=1.0,>=Stealth,font=\small]
  % input room
  \node[spacegray] at (-2.0,2.5) {input room};
  \draw[sourceblue!70!black,very thick] (-3.6,-1.4) -- (-0.4,1.4);
  \node[sourceblue!70!black,right] at (-1.0,1.25) {row space};
  \draw[quotientorange!85!black,very thick,dashed] (-3.4,1.2) -- (-0.6,-1.2);
  \node[quotientorange!85!black,left] at (-3.35,1.15) {nullspace};
  \coordinate (O) at (-2.0,0);
  \node[fill=black,circle,inner sep=1pt] at (O) {};
  % x and its two pieces
  \coordinate (xrow) at (-1.0,0.75);   % along row space dir
  \coordinate (xnul) at (-2.55,0.75);  % along nullspace dir
  \coordinate (X) at ($(xrow)+(xnul)-(O)$);
  \draw[->,black,very thick] (O) -- (X) node[above right]{$\vec{x}$};
  \draw[->,sourceblue!70!black,thick] (O) -- (xrow) node[below right=-1pt]{$\vec{x}_{\text{row}}$};
  \draw[->,quotientorange!85!black,thick] (O) -- (xnul) node[above left=-2pt]{$\vec{x}_{\text{null}}$};
  \draw[gray!60,dashed] (X)--(xrow); \draw[gray!60,dashed] (X)--(xnul);
  % arrow A
  \draw[->,structurepurple,very thick] (0.6,0.3) -- (2.0,0.3) node[midway,above]{$\mat{A}$};
  % output room
  \node[spacegray] at (4.2,2.5) {output room};
  \draw[targetgreen!80!black,very thick] (2.6,-1.2) -- (5.6,1.4);
  \node[targetgreen!80!black,right] at (4.2,1.2) {column space};
  \coordinate (O2) at (3.4,-0.2);
  \node[fill=black,circle,inner sep=1pt] at (O2) {};
  \coordinate (AX) at (4.9,1.0);
  \draw[->,targetgreen!80!black,very thick] (O2) -- (AX) node[right]{$\mat{A}\vec{x}=\mat{A}\vec{x}_{\text{row}}$};
\end{tikzpicture}
\caption{The splitting of Example~\ref{ch15:splitting}. The input $\vec{x}$ resolves into a
row-space part $\vec{x}_{\text{row}}$ and a nullspace part $\vec{x}_{\text{null}}$. Applying
$\mat{A}$ discards the nullspace part entirely, so the output depends only on $\vec{x}_{\text{row}}$
and lands in the column space.}
\label{ch15:split-arrows}
\end{figure}

\section{Reading all four bases off one elimination}

Let us consolidate the recipe, because it is genuinely efficient: a single row reduction hands you
everything.

\begin{formalbox}
\textbf{The one-elimination recipe.} Reduce $\mat{A}$ (tracking operations in $\mat{E}$ so that
$\mat{E}\mat{A}=\mat{R}$). Let the pivots sit in $r$ columns.
\begin{itemize}
  \item \textbf{Column space} ($\dim r$, in $\R^m$): the pivot columns of the \emph{original}
  $\mat{A}$.
  \item \textbf{Row space} ($\dim r$, in $\R^n$): the nonzero rows of $\mat{R}$.
  \item \textbf{Nullspace} ($\dim n-r$, in $\R^n$): the special solutions, one per free column.
  \item \textbf{Left nullspace} ($\dim m-r$, in $\R^m$): the rows of $\mat{E}$ that produced the
  zero rows of $\mat{R}$.
\end{itemize}
\end{formalbox}

\begin{labbox}
Try the full recipe on
\[
  \mat{A} = \begin{bmatrix} 1 & 2 & 1 & 1\\ 2 & 4 & 3 & 4\\ 1 & 2 & 2 & 3 \end{bmatrix}
  \quad (m=3,\ n=4).
\]
Eliminate to $\mat{R}=\begin{bmatrix}1&2&0&-1\\0&0&1&2\\0&0&0&0\end{bmatrix}$ (pivots in columns 1
and 3, so $r=2$). Now read off all four:
\begin{itemize}
  \item \textbf{Column space}: columns 1 and 3 of $\mat{A}$, namely $\colvec{1\\2\\1}$ and
  $\colvec{1\\3\\2}$. ($\dim 2$, in $\R^3$.)
  \item \textbf{Row space}: $\colvec{1\\2\\0\\-1}$ and $\colvec{0\\0\\1\\2}$. ($\dim 2$, in
  $\R^4$.)
  \item \textbf{Nullspace}: free columns are 2 and 4. Setting $x_2=1$ (rest free $=0$) gives
  $\colvec{-2\\1\\0\\0}$; setting $x_4=1$ gives $\colvec{1\\0\\-2\\1}$. ($\dim 2$, in $\R^4$.)
  \item \textbf{Left nullspace}: one zero row, so $\dim 1$ in $\R^3$. The combination
  $R_1-R_2+R_3$ of the original rows is zero, giving $\colvec{1\\-1\\1}$.
\end{itemize}
Sanity check the dimensions: $r=2$, $n-r=2$ (fills $\R^4$), $m-r=1$ (with $r=2$ fills $\R^3$).
\end{labbox}

\begin{trapbox}
Watch your rooms. The row space and nullspace are sets of $4$-component vectors (they live in the
\emph{input} space $\R^4$); the column space and left nullspace are $3$-component vectors (they live
in the \emph{output} space $\R^3$). A frequent error is to try to ``check orthogonality'' between
the column space and the nullspace -- but those two live in different rooms and cannot even be
dotted together. Orthogonality is always \emph{within} a room: row $\perp$ null in $\R^n$, column
$\perp$ left-null in $\R^m$.
\end{trapbox}

\section{A modern application: where data lives and what is unreachable}

\begin{applicationbox}
\textbf{Recommender systems and the column space.} Imagine a big matrix $\mat{A}$ whose rows are
users and whose columns are movies, with entry $A_{ij}$ the rating user $i$ gave movie $j$. Most
entries are missing; the engine's job is to fill them in. The working hypothesis behind systems
like the ones that won the Netflix Prize is that tastes are \emph{low rank}: a handful of hidden
factors (how much you like comedy, action, slow burns) explain most of the ratings. ``Low rank''
means the rank $r$ is small compared with $m$ and $n$.

The four subspaces tell you exactly what such a model can and cannot do.
\begin{itemize}
  \item The \textbf{column space} is the set of rating patterns the model can produce. If a real
  user's true ratings lie outside this $r$-dimensional column space, no choice of weights can
  reproduce them -- their genuinely individual taste is, mathematically, in the
  \textbf{left nullspace}, perpendicular to everything the model can represent. That orthogonal
  leftover is precisely the part the model must give up on.
  \item The \textbf{nullspace} is the model's blind spot among inputs: directions in
  user/feature space you can change without changing a single predicted rating. Regularizers and
  extra data are how engineers try to shrink or pin down that nullspace so predictions become
  determined rather than free.
\end{itemize}
The same decomposition powers image compression, topic models for text, and the embeddings inside
modern recommendation and language models: keep the $r$ directions that carry signal (the column
space), discard the rest. The rank is the budget; the four subspaces are the accounting.
\end{applicationbox}

\begin{applicationbox}
\textbf{Least squares and the left nullspace, previewed.} When $\mat{A}\vec{x}=\vec{b}$ has no
solution -- the overdetermined case at the heart of every regression -- it is precisely because
$\vec{b}$ has a piece sticking out of the column space, a piece lying in the left nullspace. The
least-squares method, which you will meet soon, throws that unreachable piece away by projecting
$\vec{b}$ onto the column space and solving for the part that \emph{is} achievable. The left
nullspace is not an exotic curiosity; it is the exact measure of ``how much your data refuses to
fit your model.''
\end{applicationbox}

\section{Worked summary example, all four at once}

\begin{example}\label{ch15:full-worked}
Let
\[
  \mat{A} = \begin{bmatrix} 1 & 3 & 2\\ 2 & 6 & 4 \end{bmatrix}
  \qquad (m=2,\ n=3).
\]
Row~2 is twice row~1, so we expect $r=1$. Eliminate: $R_2-2R_1$ gives
$\mat{R}=\begin{bmatrix}1&3&2\\0&0&0\end{bmatrix}$. One pivot, $r=1$.
\begin{itemize}
  \item \textbf{Column space} (in $\R^2$, $\dim 1$): pivot column~1 of $\mat{A}$ is
  $\colvec{1\\2}$. The column space is the line through $\colvec{1\\2}$.
  \item \textbf{Left nullspace} (in $\R^2$, $\dim m-r=1$): the operation $R_2-2R_1$ zeroed a row,
  so the combination $-2(\text{row}_1)+1(\text{row}_2)$ vanishes, i.e. $\vec{y}=\colvec{-2\\1}$
  gives $\vec{y}\T\mat{A}=\vec{0}\T$. Check entry by entry: $-2(1)+1(2)=0$, $-2(3)+1(6)=0$, and
  $-2(2)+1(4)=0$. Good. The left nullspace is the line through $\colvec{-2\\1}$, perpendicular to
  the column-space direction $\colvec{1\\2}$ (since $(-2)(1)+(1)(2)=0$). Together they fill
  $\R^2$.
  \item \textbf{Row space} (in $\R^3$, $\dim 1$): the nonzero reduced row $\colvec{1\\3\\2}$. The
  line through it.
  \item \textbf{Nullspace} (in $\R^3$, $\dim n-r=2$): free columns 2, 3. Special solutions
  $\colvec{-3\\1\\0}$ (set $x_2=1$) and $\colvec{-2\\0\\1}$ (set $x_3=1$). A plane. Both are
  orthogonal to the row space direction $\colvec{1\\3\\2}$: $-3+3+0=0$ and $-2+0+2=0$. The line
  (row space) plus the plane (nullspace) fill $\R^3$.
  \end{itemize}
Every dimension checks: $\R^2 = 1 + 1$, $\R^3 = 1 + 2$. A rank-one matrix: it listens to a single
input direction and writes to a single output direction, ignoring everything else.
\end{example}

\begin{figure}[H]
\centering
\begin{tikzpicture}[scale=1.0,>=Stealth]
  % R^2 output room
  \begin{scope}
    \draw[->,spacegray] (-2.2,0)--(2.2,0);
    \draw[->,spacegray] (0,-2.2)--(0,2.2);
    \node[spacegray] at (0,2.5) {output room $\R^2$};
    % column space line through (1,2)
    \draw[targetgreen!80!black,very thick] (-1,-2)--(1,2);
    \draw[->,targetgreen!80!black,very thick] (0,0)--(1,2);
    \node[targetgreen!80!black,right] at (0.7,1.7){$\col$};
    % left null line through (-2,1)
    \draw[structurepurple,very thick,dashed] (2,-1)--(-2,1);
    \draw[->,structurepurple,very thick] (0,0)--(-1.4,0.7);
    \node[structurepurple,left] at (-1.2,0.95){$\nullsp(\mat{A}\T)$};
  \end{scope}
\end{tikzpicture}
\caption{The output room $\R^2$ for the rank-one matrix of
Example~\ref{ch15:full-worked}: the column space (green, the line through
$\protect\colvec{1\\2}$) and the left nullspace (purple dashed, the perpendicular line through
$\protect\colvec{-2\\1}$). They meet at a right angle and together fill the plane.}
\label{ch15:rank1-room}
\end{figure}

\begin{historybox}
The four-subspace framing and the name ``Fundamental Theorem of Linear Algebra'' for the
orthogonality-plus-dimensions package were popularized by Gilbert Strang, partly in a 1993
\emph{American Mathematical Monthly} article written precisely to argue that this picture, not the
determinant, deserves to be the centerpiece of a first course. The mathematics is older -- the
rank-and-orthogonality story threads back through Sylvester, Frobenius, and Fredholm -- but the
\emph{pedagogy} of drawing two rooms with two perpendicular subspaces apiece and a clean
isomorphism between them is Strang's gift to generations of students.
\end{historybox}

\section{What we now know}

\begin{focusbox}
The four fundamental subspaces of an $m\times n$ matrix $\mat{A}$ of rank $r$:
\begin{center}
\begin{tabular}{lll}
\toprule
Subspace & Lives in & Dimension \\
\midrule
Column space $\col(\mat{A})$ & $\R^m$ & $r$ \\
Left nullspace $\nullsp(\mat{A}\T)$ & $\R^m$ & $m-r$ \\
Row space $\row(\mat{A})$ & $\R^n$ & $r$ \\
Nullspace $\nullsp(\mat{A})$ & $\R^n$ & $n-r$ \\
\bottomrule
\end{tabular}
\end{center}
In $\R^n$: row space $\perp$ nullspace, dimensions add to $n$.
In $\R^m$: column space $\perp$ left nullspace, dimensions add to $m$.
And $\mat{A}$ maps the row space one-to-one onto the column space, collapsing the nullspace to
$\vec{0}$.
\end{focusbox}

In the next chapter we give the rank itself center stage and state the
Fundamental Theorem of Linear Algebra in full -- you have already done most of the work.

\section{Exercises}

\begin{exercise}\label{ch15:ex-basic}
For $\mat{A}=\begin{bmatrix}1&2&0\\0&0&1\end{bmatrix}$, find a basis and the dimension of each of
the four fundamental subspaces. State which room ($\R^2$ or $\R^3$) each lives in.
\end{exercise}

\begin{exercise}\label{ch15:ex-special}
Let $\mat{A}=\begin{bmatrix}1&1&2&1\\2&2&4&3\\1&1&2&0\end{bmatrix}$. Reduce it, identify the
pivot and free columns, and write down all the special solutions (a basis for the nullspace).
What is the nullity?
\end{exercise}

\begin{exercise}\label{ch15:ex-colspace-trap}
Give a $2\times2$ matrix whose row-reduced form has a pivot in column~1, yet whose column space is
\emph{not} the $x$-axis. Use it to explain in one sentence why the column-space basis must be read
from the original matrix, not the reduced one.
\end{exercise}

\begin{exercise}\label{ch15:ex-leftnull}
For $\mat{A}=\begin{bmatrix}1&2\\2&4\\3&6\end{bmatrix}$, find the left nullspace directly by
solving $\mat{A}\T\vec{y}=\vec{0}$. Confirm its dimension equals $m-r$ and that every basis vector
is orthogonal to the column space.
\end{exercise}

\begin{exercise}\label{ch15:ex-dimcheck}
A matrix $\mat{A}$ is $5\times 7$ with rank $4$. State the dimension of each of the four
fundamental subspaces and the room each lives in. How many independent special solutions does
$\mat{A}\vec{x}=\vec{0}$ have?
\end{exercise}

\begin{exercise}\label{ch15:ex-splitting}
For $\mat{A}=\begin{bmatrix}1&1\\1&1\end{bmatrix}$ and input $\vec{x}=\colvec{3\\1}$, split
$\vec{x}$ into its row-space and nullspace pieces and verify that $\mat{A}\vec{x}=\mat{A}\vec{x}_
{\text{row}}$, the content of Theorem~\ref{ch15:iso-thm}.
\end{exercise}

\begin{exercise}\label{ch15:ex-solvability}
Use the left nullspace to decide solvability. For
$\mat{A}=\begin{bmatrix}1&2\\2&4\\0&1\end{bmatrix}$, a vector $\vec{b}$ makes
$\mat{A}\vec{x}=\vec{b}$ solvable iff $\vec{b}$ is orthogonal to the left nullspace. Find that
orthogonality condition explicitly and check it on $\vec{b}=\colvec{1\\2\\0}$ and
$\vec{b}=\colvec{1\\1\\1}$.
\end{exercise}

\begin{exercise}\label{ch15:ex-rowcol-iso}
Let $\mat{A}=\begin{bmatrix}2&0\\0&3\\0&0\end{bmatrix}$. Identify all four subspaces, then verify
by hand that $\mat{A}$ carries the row-space basis to the column-space basis in a one-to-one onto
fashion (Theorem~\ref{ch15:iso-thm}). What are the two stretch factors?
\end{exercise}

\begin{exercise}\label{ch15:ex-proof}
\emph{(Proof.)} Prove that for any $m\times n$ matrix $\mat{A}$, a vector in both the row space and
the nullspace must be the zero vector. Then use this to prove that the only vector in all four
subspaces simultaneously is $\vec{0}$ (after accounting for the rooms they live in).
\end{exercise}

\begin{exercise}\label{ch15:ex-transpose}
\emph{(Symmetry challenge.)} Explain why $\rank(\mat{A})=\rank(\mat{A}\T)$ guarantees that the
column space and row space always have the \emph{same} dimension, even when $m\ne n$ so they live
in rooms of different sizes. Illustrate with a $2\times 5$ matrix of rank $2$: what are the four
dimensions?
\end{exercise}

\begin{exercise}\label{ch15:ex-application}
\emph{(Modeling.)} A data matrix $\mat{A}$ is $1000\times 50$ (1000 samples, 50 features) and turns
out to have rank $r=8$. In plain language, what do the column space, nullspace, row space, and left
nullspace each mean for this dataset? Which subspace measures ``directions in feature space that do
not affect any output,'' and which measures ``output patterns the data can never produce''?
\end{exercise}
